\begin{proof}[证明 (精确有理包络)]
使用 Machin 恒等式
$\pi=16\arctan(1/5)-4\arctan(1/239)$ 及交错部分和得到 $\pi$ 的有理上下界.
对有理 $x$ 用 $\sin x$, $\cos x$ 的 Taylor 多项式及
$|R_N(x)|\le |x|^{N+1}/(N+1)!$ 作包络; 除法前检查分母区间不含零.
在两个精确有理端点
\begin{align*}
 a_L&=2.27651323902643307501655540074525185589672,\\
 a_H&=2.27651323902643307501655540074525185589673
\end{align*}
分别证 $G(a_L)>0$, $G(a_H)<0$ (此处 $G$ 为 T2 的一变量符号链函数).
由 T2 的唯一性得 $a_G\in(a_L,a_H)$. 使用严格递增的
$u(a)=a/[2(a-\tan a)]$ 传播区间时, 取下端像的下界和上端像的上界,
再对 $(a_G^2-\pi^2/4)/u(a_G)^2$ 使用精确区间四则运算.
所得更窄包络严格位于本定理的两个显示区间内; 比较 $25$ 和 $3\pi^2$
也在有理算术中完成. 十进制显示按向下/向上舍入, 不参与真值判定.
第八轮证书及输入身份见 §\ref{sec:cac}.
旧脚本 05 的两端均取根下界规则不能保证像包含, 不再作为本定理证书.
\end{proof}

